% ==================================================================
% id:       prob-013
% topic:    מבוא להסתברות
% subtopic: עקרון ההכלה וההפרדה
% solution: full
% ==================================================================
דני מתכונן לשני מבחנים, בחשבון ובאנגלית. להערכת המורים שלו, ההסתברות שיעבור את המבחן בחשבון היא $0.7$, ההסתברות שיעבור את המבחן באנגלית היא $0.4$, וההסתברות שיעבור לפחות אחד מהם היא $0.9$.
\begin{enumerate}
\item מה ההסתברות שדני ייכשל לפחות במבחן אחד?
\begin{solbox}
נסמן $M$ = „עובר בחשבון”, $E$ = „עובר באנגלית”. לפי הכלה והפרדה, $\PP(M\cap E)=0.7+0.4-0.9=0.2$. „נכשל לפחות באחד” הוא המשלים של „עובר את שניהם”: $1-0.2=\ans{0.8}$.
\end{solbox}
\item מה ההסתברות שדני ייכשל בשני המבחנים?
\begin{solbox}
זה המשלים של „עובר לפחות אחד”: $1-\PP(M\cup E)=1-0.9=\ans{0.1}$.
\end{solbox}
\item מה ההסתברות שדני יעבור בדיוק מבחן אחד?
\begin{solbox}
$\PP(M\cup E)-\PP(M\cap E)=0.9-0.2=\ans{0.7}$.
\end{solbox}
\end{enumerate}
